\subsection{Regular university admissions in Chile}\label{app:paes}
The public PAES applications list each student's program choices in order, their program-specific weighted score and their official status. Our capacity equals the sum of first- and second-semester regular seats and announced overbooking seats. We keep regular applications coded selected, waitlisted or passed over following a higher selection (24, 25, 26). We rank by weighted score and break score ties by a fixed seeded lottery. In a check that omits all test-year status codes when selecting candidate edges, keeping valid nonduplicate reported programs with a strictly positive score gives exactly the same applicant--program pairs, order and scores in 2026. We never distinguish selected from waitlisted when producing a 2026 display. Program-specific selection rank and individual identifiers do not enter the predictor. The public input omits five applicants in 2026 for confidentiality. A recorded offer is an invitation to enroll. Subsequent enrollment and waitlist movements are different outcomes.

The 2024 and 2025 cohorts have 188,707 and 189,982 regular applicants. For the 2025 model, we use each program's previous minimum regular offer score and reported demand per seat from 2024. Missing program cutoffs receive the median available cutoff; missing congestion receives the median across prior programs. Ridge logistic regression with a fixed feature basis uses log submitted rank, log list length, score, cutoff distance, capacity and interactions. Its label is the matched edge under full-list applicant-proposing DA (zero for every other edge). The year-ahead run applies the same coefficients to 2026 features using 2025 historical summaries. A second control keeps programs where the weighted score clears last year's program cutoff, in applicant rank order, then fills any remaining slots from the other programs. All displayed programs are proposed to in original preference order. These rules could be evaluated after applications are drafted. They are not search tools for programs the student never considered.

\begin{table}[ht]\centering\small
\caption{Chilean university menus. The first numeric column measures recovery among 156,248 recorded regular recipients. The second uses 158,697 full-list DA recipients. The third uses all 189,399 applicants. The last two columns show the placement change against uniform for caps of one and ten choices. Uniform coverage is analytic; placement averages 64 draws.}
\label{tab:paes}\begin{tabular}{@{}lrrrrr@{}}\toprule
 &\multicolumn{3}{c}{Quota five, percent}&\multicolumn{2}{c}{Placement change (pp)}\\
\cmidrule(lr){2-4}\cmidrule(l){5-6}
Rule & Recorded offer & DA partner & DA placed & Quota one & Quota ten\\\midrule
Uniform & 74.22 & 73.83 & 82.27 & --- & ---\\
First reported choices & 91.00 & 89.98 & 79.20 & $-9.68$ & $-0.87$\\
Historical cutoff & 95.90 & 96.35 & 83.74 & $+9.42$ & $+0.17$\\
Fitted classifier & 98.05 & 97.64 & 83.19 & $+8.16$ & $+0.05$\\\bottomrule
\end{tabular}\end{table}

At quota three the classifier placed 155,593 students, versus 157,932 under the cutoff rule and 150,804 on average under uniform. At five, 96,310 applicants had enough options for pruning to bind. The standard error of the mean uniform placement count was 7.8 applicants (0.0041 percentage points). The deterministic policies have no corresponding display-sampling error. Removing overbooking reduces full-list agreement with recorded offers to 105,097 of 156,248 recipients (67.26\%). The fitted rule still places more applicants than uniform and fewer than the historical-cutoff control at five choices (139,039 versus 138,448 and 139,799). Two independently written DA algorithms, sequential and simultaneous, gave identical full-market assignments. Every evaluated display passed capacity and no-blocking-pair checks. A max-flow calculation verified each displayed capacity with an equal-valued cut.
\paragraph{Prediction from recorded offers.}
As a follow-up we fitted equal-size, row-permutation-equivariant networks to the actual 2025 regular offer or the outside option. Both networks receive program selection score, score relative to the previous cutoff, previous demand per seat, current capacity, and a missing-history indicator. One also receives reported preference rank. Options are permuted independently in every applicant row, lifting any dependence on column position. Each setting uses three H200 fits, a fixed 12 epochs and average predicted probabilities. There are no 2026 labels in the inference files. Both models are evaluated at the same four quotas in the same 2026 cohort as the primary analysis. The outcome-blind construction found the same 1,328,381 eligible pairs.

With five programs, the rank-blind model retained 144,401 of the recorded regular offers (92.42\%) and placed 158,800 applicants. Giving the same architecture the submitted rank retained 155,657 offers (99.62\%) and placed 158,533. This is a gain of 11,256 covered historical offers and a net loss of 267 counterfactual placements. Relative to full-list DA, the rank-informed model caused fewer members of the original placed set to lose a place (364 versus 1,045). It also placed far fewer previously unmatched people (200 versus 1,148). The additional 681 retained recipients and 948 fewer entrants sum to the 267-placement loss. This distinction is important when the service both predicts historical offers and evaluates policy performance on previous recipients.

\subsection{Validating against observed university registrations}\label{app:enroll}
We used the ministry's separate PAES registration file, with applicant pseudonyms used only to link records locally \citep{paes}. It contains 146,873 unique students for 2026. Of these, 10,122 fall outside our regular applicant cohort. Of the 189,399 regular applicants, 136,751 have a link and 52,648 have no registration record in this release. The coverage denominator is the 136,751 linked applicants, including 5,271 whose destination is absent from the retained regular graph. The file covers several admission routes. No record here need not mean that the applicant did not attend any institution.

\begin{table}[ht]\centering\small
\caption{Five-program lists produced by the two offer-trained networks. Historical regular offer describes the program in which a registered applicant was initially invited to enroll. Gains and losses count individual enrolled destinations added or removed by supplying the network with submitted rank.}
\label{tab:enroll}\begin{tabular}{@{}lrrrrr@{}}\toprule
Historical regular offer & Students & Blind coverage & Informed coverage & Gained & Lost\\\midrule
Registered there & 116,638 & 107,773 & 116,284 & 8,580 & 69\\
Registered elsewhere & 10,264 & 5,758 & 7,063 & 1,446 & 141\\
No initial regular offer & 9,849 & 5,367 & 6,321 & 976 & 22\\\midrule
All linked applicants & 136,751 & 118,898 & 129,668 & 11,002 & 232\\\bottomrule
\end{tabular}\end{table}
\paragraph{Changing what the network learns.}
The offer comparison prompted a paired follow-up on the same cohort. We replaced only the target in the 2025 training data, retaining each applicant's enrolled program when it appeared on the regular applications, or the outside class otherwise. This changes 45,450 of 189,982 training labels. We used the same row masks, features, shuffling, two architectures, seeds, optimizer, epochs and 2026 test features. The outside class describes missing on-list enrollment and includes registration through other routes. We selected no epoch or scoring rule using the test results.

All six H200 runs completed. With enrollment as the training target, the rank-blind and rank-informed menus covered 119,414 and 130,305 of the 136,751 linked registrations. The informed menus thus covered 10,891 more actual destinations. Replayed DA placed 158,498 and 158,357 applicants, respectively, a loss of 141. The informed model kept 612 more original full-list recipients placed (1,239 versus 627 losses) but brought in 753 fewer previously unmatched applicants (1,040 versus 287 entrants). Training on registrations therefore preserved the within-architecture sign pair on this cohort. It also reveals a distinction between objectives: versus the offer-trained informed model, the enrollment-trained informed model covered 637 more registered destinations, covered 79 fewer initial offers, and placed 176 fewer applicants. The individual-level records and predictions stay out of the shared archive.


